% =============================================================================
%  Programação genérica
%  Cobre a Etapa 3 da trilha: templates, concepts + compile_time e lambdas
%  (Tier 2 do ROADMAP, inseridos aqui pelo pré-requisito).
% =============================================================================

% =============================================================================
\section{Programação genérica}
\label{sec:generica}
% =============================================================================

O C++ tem duas linguagens dentro de si. Uma roda quando o programa roda; a outra
roda quando o programa \emph{compila}, e é sobre ela que este capítulo trata.

A ordem interna importa mais aqui do que em qualquer outro capítulo.
\projeto{templates} vem antes de \projeto{concepts} porque \emph{concepts} é a
resposta do C++20 a um problema que só se entende depois de ter sofrido com
\gls{sfinae} --- estudar na ordem inversa inverte a história e transforma uma
solução elegante em sintaxe arbitrária. \projeto{compile\_time} vem depois
porque \texttt{constexpr} e \emph{templates} são as duas metades da mesma
capacidade de calcular antes de rodar. E \projeto{lambdas} fecha o capítulo
porque a discussão central dele --- \classe{std::function} contra parâmetro de
\emph{template} --- só faz sentido para quem já sabe o que um parâmetro de
\emph{template} custa.

% =============================================================================
\subsection{Templates}
\label{sec:templates}
% =============================================================================

Um \emph{template} não é uma função nem uma classe: é uma \textbf{receita} a
partir da qual o compilador gera funções e classes. A consequência que governa
tudo o mais é que cada instanciação gera código \emph{separado}, especializado
para aquele tipo --- daí virem juntos o desempenho (nada é resolvido em tempo de
execução) e os dois custos clássicos: binário maior e mensagens de erro
proporcionais à profundidade da instanciação.

\subsubsection{As formas básicas}

O subprojeto \projeto{templates} percorre a progressão inteira. As duas formas
de partida são o \emph{template} de função, cujo parâmetro é deduzido do
argumento, e o de classe, que até o C++17 exigia o argumento escrito à mão:

\begin{lstlisting}[style=bt, language=C++,
    caption={Dedução na função; CTAD, a partir do C++17, também na classe.}]
template<typename T>
T max(T a, T b) { return (a > b) ? a : b; }

max(3, 7);          // T = int, deduzido
max(3.14, 2.71);    // T = double, deduzido
max<std::string>("abc", "xyz");   // T explícito
\end{lstlisting}

E há uma família inteira de formas que vale conhecer antes de precisar delas
(\cref{tab:formas-template}), porque cada uma resolve um problema que as outras
não resolvem.

\begin{table}[htbp]
\centering
\caption{As formas de \emph{template} da linguagem.}
\label{tab:formas-template}
\footnotesize
\begin{tabular}{@{}ll@{}}
\toprule
\textbf{Forma} & \textbf{O que declara} \\
\midrule
\texttt{template <typename T> T maior(T, T);}            & \emph{template} de função \\
\texttt{template <typename T> class Pilha \{ \};}        & \emph{template} de classe \\
\texttt{template <typename T, std::size\_t N> class Buf;}& NTTP: um \emph{valor} como parâmetro \\
\texttt{template <typename T = int> struct V \{ \};}     & argumento padrão \\
\texttt{template <> struct X<int> \{ \};}                & especialização \textbf{total} \\
\texttt{template <typename T> struct X<T*> \{ \};}       & especialização \textbf{parcial} (só classes) \\
\texttt{template <typename T> constexpr bool eh\_v = …;} & \emph{template} de variável \\
\texttt{template <typename T> using Par = std::pair<T,T>;}& \emph{template} de \emph{alias} \\
\bottomrule
\end{tabular}
\end{table}

\begin{armadilha}[nomes dependentes exigem \texttt{typename}]
Dentro de um \emph{template}, \texttt{T::value\_type} é ambíguo para o
compilador: pode ser um tipo ou um membro estático. Ele supõe \emph{valor} por
padrão, e a correção é escrever \texttt{typename T::value\_type}. O sintoma é
uma mensagem de erro que reclama de sintaxe numa linha que parece
sintaticamente perfeita. A mesma história vale para \texttt{template} antes de
um membro \emph{template} dependente.
\end{armadilha}

\begin{armadilha}[especialização parcial não existe para funções]
Só \emph{templates de classe} podem ser parcialmente especializados. Para
funções, a ferramenta correta é a resolução de sobrecarga --- ou, hoje, um
\emph{concept} (\cref{sec:concepts}). Tentar especializar parcialmente uma
função leva a um erro obscuro ou, pior, a uma sobrecarga que compila e nunca é
escolhida.
\end{armadilha}

\subsubsection{Variádicos e \emph{fold expressions}}

Antes do C++11, aceitar um número variável de argumentos tipados exigia escrever
uma sobrecarga para cada aridade. Os \emph{parameter packs} substituem toda essa
família por um \emph{template}, e as \emph{fold expressions} do C++17
substituem a recursão que os processava:

\begin{lstlisting}[style=bt, language=C++,
    caption={O mesmo problema, em três gerações da linguagem.}]
// C++03: uma sobrecarga por aridade — e nenhuma verificação de tipo na variádica C
void imprimir() {}                                    // C++11: caso base

template<typename First, typename... Rest>            // C++11: recursão
void imprimir(First first, Rest... rest) {
    fmt::print("{} ", first);
    imprimir(rest...);
}

template<typename... Ts>                              // C++17: fold expression
auto soma(Ts... vals) { return (vals + ...); }
\end{lstlisting}

A \emph{fold} tem quatro formas, e a diferença entre elas é a associatividade ---
que importa quando o operador não é associativo (\cref{tab:folds}).

\begin{table}[htbp]
\centering
\caption{As formas de \emph{fold expression}.}
\label{tab:folds}
\footnotesize
\begin{tabular}{@{}ll@{}}
\toprule
\textbf{Escrito} & \textbf{Expande para} \\
\midrule
\texttt{(args + ...)}    & \texttt{a1 + (a2 + (a3 + a4))} \\
\texttt{(... + args)}    & \texttt{((a1 + a2) + a3) + a4} \\
\texttt{((f(args)), ...)}& \texttt{f(a1), f(a2), f(a3), …} \\
\texttt{((vals > 0) \&\& ...)} & conjunção de todos \\
\bottomrule
\end{tabular}
\end{table}

\subsubsection{SFINAE, e por que ele existiu}

\gls{sfinae} é uma regra de resolução de sobrecarga: se substituir os parâmetros
de \emph{template} produzir um tipo inválido \emph{na assinatura}, o candidato é
apenas descartado, sem erro. A regra existe para permitir sobrecargas que se
excluem mutuamente conforme uma propriedade do tipo:

\begin{lstlisting}[style=bt, language=C++,
    caption={Duas sobrecargas mutuamente exclusivas, à moda pré-C++20.}]
template<typename T>
typename std::enable_if<std::is_integral<T>::value, T>::type
dobro(T x) { return x * 2; }

template<typename T>
typename std::enable_if<std::is_floating_point<T>::value, T>::type
dobro(T x) { return x * 2.0; }
\end{lstlisting}

E o \emph{detection idiom} é a aplicação mais elaborada da mesma regra: descobrir
se um tipo tem determinado membro, sem exigir que ele o tenha.

\begin{lstlisting}[style=bt, language=C++,
    caption={O \emph{detection idiom}: especialização parcial que só é válida se
             \texttt{.size()} existir.}]
template<typename T, typename = void>
struct tem_size : std::false_type {};

template<typename T>
struct tem_size<T, std::void_t<decltype(std::declval<T>().size())>>
    : std::true_type {};

static_assert(tem_size<std::vector<int>>::value);
static_assert(!tem_size<int>::value);
\end{lstlisting}

\begin{porque}[por que ninguém escreve isso de bom grado]
As três linhas acima expressam ``\emph{T} tem \texttt{size()}'' --- uma ideia que
cabe em cinco palavras --- por meio de uma especialização parcial cujo segundo
parâmetro existe apenas para ser descartado. É ilegível de propósito nenhum:
ninguém projetou o SFINAE para isso, ele apenas \emph{acabou} servindo. Quando
falha, a mensagem cita o interior da instanciação em vez do erro do chamador. É
esse acúmulo que os \emph{concepts} vieram desfazer.
\end{porque}

\subsubsection{\emph{Type traits} e \texttt{if constexpr}}

Os \emph{type traits} são a interface de introspecção: perguntas
(\texttt{is\_integral\_v}, \texttt{is\_polymorphic\_v},
\texttt{is\_nothrow\_move\_constructible\_v}) e transformações
(\texttt{remove\_const\_t}, \texttt{decay\_t}, \texttt{add\_pointer\_t}).
\texttt{decay\_t} merece nota: é exatamente o que \texttt{auto} faz ao deduzir
por valor --- \texttt{decay\_t<const int\&>} é \texttt{int}, e
\texttt{decay\_t<int[5]>} é \texttt{int*}.

O \texttt{if constexpr} do C++17 substituiu boa parte do SFINAE em código novo,
porque escolhe o galho \emph{descartando} o outro --- que nem precisa ser válido
para aquele tipo:

\begin{lstlisting}[style=bt, language=C++,
    caption={Um corpo, três comportamentos, nenhuma sobrecarga.}]
template<typename T>
void processar(T val) {
    if constexpr (std::is_integral_v<T>) {
        fmt::print("inteiro: {} (dobro: {})\n", val, val * 2);
    } else if constexpr (std::is_floating_point_v<T>) {
        fmt::print("float: {:.4f}\n", val);
    } else {
        fmt::print("outro tipo\n");
    }
}
\end{lstlisting}

\subsubsection{CRTP: polimorfismo sem \emph{vtable}}

O \gls{crtp} é a classe derivada passando \emph{a si mesma} como argumento de
\emph{template} da base. Isso dá à base o tipo exato do derivado em tempo de
compilação, e portanto o direito de chamá-lo por \texttt{static\_cast} em vez de
por \gls{vtable}:

\begin{lstlisting}[style=bt, language=C++,
    caption={A base chamando a derivada, em
             \texttt{src/templates/demo\_crtp.hpp}.}]
    // A base chama métodos da derivada via static_cast — zero overhead
    double area() const {
        return static_cast<const Derived*>(this)->area_impl();
    }

    double perimetro() const {
        return static_cast<const Derived*>(this)->perimetro_impl();
    }
\end{lstlisting}

A aplicação que mais aparece, porém, não é substituir \texttt{virtual}: é o
\emph{mixin} que deriva operações a partir de poucas primitivas. O usuário
escreve \texttt{==} e \texttt{<}; os outros quatro operadores chegam de graça:

\begin{lstlisting}[style=bt, language=C++,
    caption={Um \emph{mixin} de comparação, em
             \texttt{src/templates/demo\_crtp.hpp}.}]
struct Comparavel {
    bool operator!=(const Derived& outro) const {
        return !(static_cast<const Derived&>(*this) == outro);
    }
    bool operator>(const Derived& outro) const {
        return outro < static_cast<const Derived&>(*this);
    }
    bool operator<=(const Derived& outro) const {
        return !(static_cast<const Derived&>(*this) > outro);
    }
    bool operator>=(const Derived& outro) const {
        return !(static_cast<const Derived&>(*this) < outro);
    }
};
\end{lstlisting}

\begin{regra}[em C++20, esse \emph{mixin} específico virou uma linha]
O \texttt{operator<=>} (\cref{sec:operators}) deriva os seis operadores de
comparação a partir de um \texttt{= default}. O CRTP continua valendo para
\emph{outras} famílias de operações derivadas --- é o mecanismo de
\texttt{std::ranges::view\_interface} (\cref{sec:ranges}) --- mas comparação já
não é motivo para escrevê-lo.
\end{regra}

\begin{table}[htbp]
\centering
\caption{CRTP e \texttt{virtual} resolvem problemas parecidos com trocas opostas.}
\label{tab:crtp-virtual}
\footnotesize
\begin{tabular}{@{}lll@{}}
\toprule
\textbf{Aspecto} & \texttt{virtual} & \textbf{CRTP} \\
\midrule
resolução                 & execução        & compilação \\
custo                     & vtable + vptr   & zero \\
container polimórfico     & \texttt{vector<Base*>} & \textbf{não} (tipos diferentes) \\
código gerado             & um por hierarquia & um por tipo \\
mensagens de erro         & legíveis        & verbosas \\
\bottomrule
\end{tabular}
\end{table}

A terceira linha é a decisiva. CRTP não tem base comum em tempo de execução:
\texttt{Circulo} e \texttt{Quadrado} derivados de \texttt{FormaBase<Circulo>} e
\texttt{FormaBase<Quadrado>} não têm nenhum tipo em comum, e portanto não podem
ir para o mesmo container. Se o programa precisa de uma coleção heterogênea,
CRTP não é opção --- e é aí que o \cref{sec:ecs} propõe uma terceira via.

% =============================================================================
\subsection{Concepts}
\label{sec:concepts}
% =============================================================================

O subprojeto \projeto{concepts} começa pela comparação que justifica a
existência da funcionalidade:

\begin{lstlisting}[style=bt, language=C++,
    caption={A mesma restrição, antes e depois do C++20.}]
template<typename T>
typename std::enable_if<std::is_integral<T>::value, T>::type
soma(T a, T b) { return a + b; }     // C++17

template<std::integral T>
T soma(T a, T b) { return a + b; }   // C++20
\end{lstlisting}

Um \gls{concept} é um predicado booleano sobre tipos, avaliado em compilação. A
diferença prática em relação ao \texttt{static\_assert} é \emph{quando} ele age,
e essa diferença muda o resultado do programa.

\subsubsection{As quatro formas de escrever a mesma restrição}

\begin{lstlisting}[style=bt, language=C++,
    caption={Quatro sintaxes equivalentes.}]
template <Somavel T>            T f(T, T);              // (a) no lugar do typename
template <typename T> requires Somavel<T> T f(T, T);    // (b) cláusula requires
template <typename T> T f(T, T) requires Somavel<T>;    // (c) requires no fim
auto f(Somavel auto a, Somavel auto b);                 // (d) abreviada
\end{lstlisting}

A forma (d) é a mais curta e a mais enganosa: cada \texttt{auto} restringido
introduz um parâmetro de \emph{template} \emph{independente}, de modo que
\texttt{f(1, 2.0)} compila --- os dois \texttt{Somavel} podem ser tipos
diferentes. Quando os tipos precisam coincidir, use a forma (a).

\subsubsection{As quatro espécies de \emph{requirement}}

Dentro de um \texttt{requires}, quatro coisas diferentes podem ser exigidas, e o
subprojeto tem um exemplo de cada.

\paragraph{\emph{Simple requirement}} --- a expressão precisa compilar; nada se
diz sobre o tipo do resultado:

\begin{lstlisting}[style=bt, language=C++,
    caption={\emph{Simple requirements}, em
             \texttt{src/concepts/demo\_requires.hpp}.}]
// Concept: tipo que pode ser somado com si mesmo
template <typename T>
concept Somavel = requires(T a, T b) {
    a + b;        // simple requirement: a + b precisa compilar
    a - b;        // a - b também precisa compilar
    a += b;       // a += b também
};
\end{lstlisting}

\paragraph{\emph{Type requirement}} --- um tipo aninhado precisa existir. É a
forma de exigir o contrato de container discutido na \cref{sec:stl}:

\begin{lstlisting}[style=bt, language=C++,
    caption={\emph{Type requirements} combinados com expressões, em
             \texttt{src/concepts/demo\_requires.hpp}.}]
template <typename T>
concept Container = requires {
    typename T::value_type;        // type requirement: value_type existe
    typename T::iterator;          // type requirement: iterator existe
    typename T::const_iterator;    // type requirement: const_iterator existe
} && requires(T c) {
    c.begin();                     // begin() existe
    c.end();                       // end() existe
    c.size();                      // size() existe
};
\end{lstlisting}

\paragraph{\emph{Compound requirement}} --- a expressão precisa compilar
\emph{e} o resultado precisa satisfazer outro \emph{concept}:
\texttt{\{ a + b \} -> std::same\_as<T>;}. É a forma que verifica tipo de
retorno, e a que torna \emph{concepts} estritamente mais expressivos que
\texttt{enable\_if} sobre \emph{traits}.

\paragraph{\emph{Nested requirement}} --- um \texttt{requires} dentro do
\texttt{requires}, que aplica outro \emph{concept} ou uma condição booleana:

\begin{lstlisting}[style=bt, language=C++,
    caption={\emph{Nested requirements}, em
             \texttt{src/concepts/demo\_requires.hpp}.}]
// Concept: tipo numérico de precisão adequada (≥ 4 bytes)
template <typename T>
concept NumericoPreciso = requires {
    requires std::floating_point<T>;          // nested: sub-concept aplicado
    requires (sizeof(T) >= sizeof(float));    // nested: constraint booleana
};
\end{lstlisting}

\begin{armadilha}[um \texttt{requires} sem \texttt{requires} é sempre verdadeiro]
A construção \texttt{requires \{ expr; \}} verifica se \texttt{expr}
\emph{compila}; ela não verifica se \texttt{expr} é \emph{verdadeira}. Escrever
\texttt{requires \{ sizeof(T) >= 4; \}} é um \emph{simple requirement} que
pergunta se a comparação compila --- e ela sempre compila. A forma correta
precisa do segundo \texttt{requires}, como acima. É o erro mais comum de quem
começa, e ele não produz nenhum diagnóstico: o \emph{concept} passa a aceitar
tudo.
\end{armadilha}

\subsubsection{\emph{Subsumption}}

Esta é a propriedade que nenhuma emulação com \texttt{enable\_if} oferecia sem
trabalho manual: quando dois \emph{templates} restritos são candidatos, o
compilador escolhe o \textbf{mais específico}, e sabe qual é o mais específico
comparando os \emph{concepts} termo a termo.

\begin{lstlisting}[style=bt, language=C++,
    caption={A escolha automática da restrição mais forte.}]
template<typename T> concept Numerico = std::is_arithmetic_v<T>;
template<typename T> concept NumericoInteiro = Numerico<T> && std::integral<T>;

template<Numerico T>        void processar(T x) { /* genérico */ }
template<NumericoInteiro T> void processar(T x) { /* inteiro  */ }

processar(3.14);   // genérico — double não é integral
processar(42);     // inteiro  — satisfaz os dois, e NumericoInteiro subsume
\end{lstlisting}

Com \texttt{enable\_if}, as duas versões seriam ambíguas para \texttt{int}, e a
desambiguação exigiria negar explicitamente a condição na versão genérica.

\begin{regra}[\emph{concept} para restringir, \texttt{static\_assert} para afirmar]
A diferença é \emph{quando} cada um age:
\begin{center}\footnotesize
\begin{tabular}{@{}lll@{}}
\toprule
\textbf{Ferramenta} & \textbf{Quando falha} & \textbf{Afeta a sobrecarga?} \\
\midrule
\texttt{static\_assert} & \textbf{depois} da escolha & não \\
\emph{concept}          & \textbf{durante} a escolha & \textbf{sim} \\
\bottomrule
\end{tabular}
\end{center}
Um \texttt{static\_assert} dentro de um \emph{template} já escolhido não deixa
outra sobrecarga assumir --- ele apenas transforma a escolha ruim num erro. O
\emph{concept} remove o candidato da disputa, e a sobrecarga não restrita vence
normalmente.
\end{regra}

\subsubsection{Onde \emph{concepts} rendem mais}

Duas aplicações do subprojeto merecem destaque porque não são sobre restringir
um argumento, e sim sobre \textbf{escrever um contrato}.

A primeira é o algoritmo genérico que declara exatamente o que exige do
\emph{range} e do predicado --- uma assinatura que documenta a si mesma e que
falha no ponto da chamada, não trinta níveis abaixo:

\begin{lstlisting}[style=bt, language=C++,
    caption={O contrato inteiro na assinatura.}]
template<std::ranges::input_range Range, typename Pred>
    requires std::predicate<Pred, std::ranges::range_value_t<Range>>
auto filtrar(const Range& range, Pred pred);
\end{lstlisting}

A segunda é o \emph{policy-based design}: um \emph{concept} descreve o que uma
política precisa oferecer, e a classe que a recebe passa a recusar em compilação
qualquer política incompleta --- com uma mensagem que nomeia o método que falta,
em vez de apontar para a linha em que ele seria chamado.

\begin{table}[htbp]
\centering
\caption{SFINAE e \emph{concepts}, ponto a ponto.}
\label{tab:sfinae-concepts}
\footnotesize
\begin{tabular}{@{}lll@{}}
\toprule
\textbf{Aspecto} & \textbf{SFINAE} & \textbf{Concepts} \\
\midrule
legibilidade         & ruim (aninhamento, \texttt{enable\_if}) & boa \\
mensagens de erro    & gigantescas                             & citam a restrição violada \\
\emph{subsumption}   & manual e verbosa                        & automática \\
verificação          & implícita, por efeito colateral         & explícita, via \texttt{requires} \\
reutilização         & difícil                                 & trivial (são nomeados) \\
\bottomrule
\end{tabular}
\end{table}

% =============================================================================
\subsection{Computação em tempo de compilação}
\label{sec:compile-time}
% =============================================================================

O subprojeto \projeto{compile\_time} trata das quatro palavras-chave que movem
trabalho da execução para a compilação --- \texttt{constexpr},
\texttt{consteval}, \texttt{constinit}, \texttt{static\_assert} --- e dos
parâmetros de \emph{template} que são valores. Compila com C++23, por causa de
\texttt{if consteval} (\cref{sec:ambiente}).

O argumento a favor não é apenas desempenho. É este:

\begin{porque}[em contexto \texttt{constexpr}, \gls{ub} vira erro de compilação]
\begin{lstlisting}[style=bt, language=C++, frame=none,
                   xleftmargin=0pt, xrightmargin=0pt]
constexpr int x = 1 / 0;      // erro de compilação, não crash em produção
constexpr int y = arr[10];    // erro se fora dos limites
\end{lstlisting}
O avaliador de expressões constantes é obrigado a recusar qualquer operação cujo
comportamento o padrão não define. Código que passa sob \texttt{static\_assert}
é código \textbf{provadamente} livre de \gls{ub} naquele caminho --- uma
garantia que nenhuma bateria de testes oferece, porque testes cobrem entradas e
isto cobre o caminho inteiro.
\end{porque}

\subsubsection{\texttt{constexpr} é permissão; \texttt{consteval} é obrigação}

A mesma função \texttt{constexpr} serve aos dois mundos: chamada num contexto
que exige constante, é avaliada na compilação; chamada com um argumento de
execução, é compilada normalmente.

\begin{table}[htbp]
\centering
\caption{As três palavras, e o que cada uma exige.}
\label{tab:constexpr-consteval}
\footnotesize
\begin{tabular}{@{}lcc@{}}
\toprule
\textbf{Palavra} & \textbf{Avaliação em compilação} & \textbf{Argumento de execução} \\
\midrule
(nenhuma)            & não                & aceita \\
\texttt{constexpr}   & \textbf{permitida} & aceita \\
\texttt{consteval}   & \textbf{obrigatória} & \textbf{erro} \\
\bottomrule
\end{tabular}
\end{table}

Forçar a avaliação em compilação, quando a função é apenas \texttt{constexpr},
se faz pelo contexto: \texttt{constexpr int a = f();} força naquela declaração;
\texttt{static\_assert(f() == x);} força e ainda verifica; e
\texttt{std::array<int, f()>} força, porque um NTTP é o contexto mais exigente
que existe.

A aplicação mais útil de \texttt{consteval} é a \textbf{validação de argumento
literal}. Como um \texttt{throw} que escapa de uma avaliação constante é erro de
compilação, o construtor abaixo recusa portas inválidas antes de o programa
existir:

\begin{lstlisting}[style=bt, language=C++,
    caption={Validação em compilação, em
             \texttt{src/compile\_time/demo\_consteval.hpp}.}]
class Porta {
public:
    consteval explicit Porta(int valor) : valor_(valor) {
        if (valor < 1 || valor > 65535) {
            // Numa avaliação em compilação, um throw que escapa é ERRO DE
            // COMPILAÇÃO. A mensagem do compilador cita esta linha.
            throw "porta fora do intervalo 1..65535";
        }
    }
    constexpr int valor() const { return valor_; }
private:
    int valor_;
};
\end{lstlisting}

É exatamente o mecanismo do \texttt{std::format} do C++20:
\texttt{std::format("\{:d\}", "texto")} falha na \textbf{compilação}, e não com
uma exceção em produção.

\begin{armadilha}[\texttt{std::is\_constant\_evaluated()} tem duas pegadinhas]
Dentro de um \texttt{if} comum, ele responde \texttt{true} sempre que a função
\emph{é} \texttt{constexpr}, mesmo quando a chamada está sendo compilada para
execução. E \texttt{if constexpr (std::is\_constant\_evaluated())} é
\textbf{sempre} verdadeiro, porque a própria condição é avaliada como constante.
O \texttt{if consteval} do C++23 existe para não ter nenhuma das duas.
\end{armadilha}

\begin{armadilha}[\texttt{constexpr std::vector} não compila]
Desde o C++20 é possível usar \texttt{std::vector} e \texttt{std::string}
\emph{durante} uma avaliação constante --- mas vale a \textbf{regra da alocação
transitória}: tudo o que for alocado precisa ser liberado antes de a avaliação
terminar, porque a memória do compilador não sobrevive até a execução. Logo, um
\texttt{vector} pode ser usado no meio do cálculo e nunca ser o resultado dele.
Para devolver dados de uma função \texttt{constexpr}, o tipo é
\texttt{std::array}.
\end{armadilha}

O terceiro membro da família resolve um problema diferente. \texttt{constinit}
exige inicialização em compilação, mas \textbf{permite} mutação depois --- o que
o torna a resposta direta ao \emph{static initialization order fiasco}: se o
compilador recusa qualquer inicialização que não caiba na compilação, o fiasco
deixa de ser um bug de execução e vira um erro de \emph{build}.

\subsubsection{\texttt{static\_assert}: onde não há substituto}

\begin{regra}[a mensagem é o produto]
Uma mensagem de \texttt{static\_assert} deve dizer \textbf{o que fazer}, não o
que falhou --- quem lê já sabe o que falhou, porque o compilador cita a
condição.
\begin{lstlisting}[style=bt, language=C++, frame=none,
                   xleftmargin=0pt, xrightmargin=0pt]
static_assert(std::is_trivially_copyable_v<T>,
    "serializar() copia bytes crus e só aceita tipos trivialmente copiáveis. "
    "Tipos com ponteiros, std::string ou destrutor próprio precisam de uma "
    "sobrecarga específica.");
\end{lstlisting}
Custa uma linha e poupa a próxima pessoa de reconstruir a intenção a partir do
código.
\end{regra}

E há um domínio em que \emph{concepts} não substituem \texttt{static\_assert},
porque \emph{concepts} restringem \emph{templates} e não dizem nada sobre
\emph{layout}: a verificação de \gls{abi}.

\begin{lstlisting}[style=bt, language=C++,
    caption={Afirmações sobre \emph{layout} que nenhum \emph{concept} expressa.}]
static_assert(sizeof(CabecalhoDeRede) == 8,
    "O cabeçalho tem 8 bytes por especificação. Se este assert falhou, o "
    "padding mudou e os pacotes serão incompatíveis com a outra ponta.");
static_assert(offsetof(CabecalhoDeRede, comprimento) == 2);
\end{lstlisting}

\subsubsection{NTTP: o valor entra no tipo}

Um parâmetro de \emph{template} que não é um tipo faz o \emph{valor} participar
da identidade do tipo. \texttt{BufferFixo<int,4>} e \texttt{BufferFixo<int,16>}
são tipos \textbf{diferentes}, e \texttt{capacidade()} não lê memória nenhuma.

A aplicação que mais rende é verificação dimensional --- a mesma ideia dos
literais de unidade da \cref{sec:operators}, agora com dimensões de matriz:

\begin{lstlisting}[style=bt, language=C++,
    caption={Erro de dimensão detectado em compilação.}]
Matriz<2,3> * Matriz<3,4> → Matriz<2,4>     // compila
Matriz<2,3> * Matriz<2,2> → NÃO COMPILA     // dimensões incompatíveis
\end{lstlisting}

Sem NTTP, as dimensões seriam campos, e o erro só apareceria em tempo de
execução --- provavelmente como leitura fora dos limites.

O C++20 permitiu NTTP de tipo classe, o que finalmente admite \emph{strings}
como parâmetro de \emph{template}. A condição é o tipo ser \emph{structural}:
todos os membros públicos, sem \texttt{private}, sem \texttt{virtual}, e os
membros também \emph{structural} --- razão pela qual \texttt{std::string} não
serve, e um \texttt{StringFixa} escrito à mão serve.

\begin{armadilha}[cada valor distinto é uma instanciação nova]
NTTP tem um custo que não aparece em nenhum \emph{benchmark} de execução:
binário maior e compilação mais lenta, proporcionalmente ao número de valores
distintos usados. Use quando os valores forem poucos, fixos e participarem da
\emph{correção} do programa. Não use quando o valor vier de configuração ou de
entrada --- e, se a intenção era apenas evitar uma alocação, o instrumento certo
é \texttt{std::span} (\cref{sec:vocabulary}).
\end{armadilha}

\subsubsection{Gerando dados na compilação}

O demo final gera uma tabela de CRC-32 e um mapa de códigos HTTP em tempo de
compilação, e depois \emph{prova} onde eles foram parar:

\begin{lstlisting}[style=bt, language=bash,
    caption={Saída de \texttt{src/compile\_time/}: as tabelas no binário.}]
$ objdump -t dist/bin/compile_time | grep -i TABELA
.rodata      0x400  TABELA_CRC32
.data.rel.ro 0x0d8  TABELA_HTTP
\end{lstlisting}

As seções diferem, e a razão é instrutiva: \texttt{TABELA\_CRC32} contém apenas
\texttt{uint32\_t} --- valores puros, que vão para \texttt{.rodata}.
\texttt{TABELA\_HTTP} contém \texttt{string\_view}, que são \textbf{ponteiros}:
o endereço só é conhecido quando o \emph{loader} mapeia o binário, e por isso
ela vai para \texttt{.data.rel.ro} --- somente-leitura \emph{após} a relocação.
Nenhuma das duas está em \texttt{.bss}, o que indicaria inicialização em tempo
de execução.

Como \texttt{std::sort} é \texttt{constexpr} desde o C++20, a tabela pode ser
escrita fora de ordem e ordenada pelo compilador --- e a busca binária em tempo
de execução passa a ser válida porque a ordenação foi \emph{provada}, não
testada:

\begin{lstlisting}[style=bt, language=C++,
    caption={A pré-condição da busca binária, verificada na compilação.}]
static_assert(std::is_sorted(TABELA_HTTP.begin(), TABELA_HTTP.end()));
\end{lstlisting}

\begin{table}[htbp]
\centering
\caption{Para poucas dezenas de entradas fixas, o array ordenado ganha nas três
         colunas.}
\label{tab:mapa-constexpr}
\footnotesize
\begin{tabular}{@{}lccc@{}}
\toprule
\textbf{Estrutura} & \texttt{constexpr} & \textbf{Aloca} & \textbf{Busca} \\
\midrule
\texttt{std::map}            & não & sim & $O(\log n)$ \\
\texttt{std::unordered\_map} & não & sim & $O(1)$ médio \\
\textbf{array ordenado}      & \textbf{sim} & \textbf{não} & $O(\log n)$ \\
\bottomrule
\end{tabular}
\end{table}

\begin{armadilha}[uma API \texttt{constexpr} robusta usa índices, não ponteiros]
O subprojeto registra uma lição que só aparece na prática:
\texttt{return p ? *p : padrao;} parece natural, mas converter um
ponteiro-para-membro-de-objeto-\texttt{constexpr} em \texttt{bool} \textbf{não é
expressão constante} em todos os modos do GCC --- o mesmo código compila
normalmente e falha sob \texttt{-fsanitize}. Aritmética de \emph{índice} é
sempre constante, e por isso é a base mais robusta.
\end{armadilha}

\begin{regra}[o custo é pago pelo compilador, e ele é lento]
O interpretador de \texttt{constexpr} do GCC é ordens de grandeza mais lento que
código nativo, e tem teto (\texttt{-fconstexpr-ops-limit}, 33,5 milhões de
passos). Não vale a pena quando a tabela é grande e o \emph{build} fica lento,
quando os dados mudam com frequência, ou quando o cálculo já era trivial na
inicialização. Vale muito quando a tabela é pequena e imutável, quando se quer
\emph{provar} uma propriedade, em ambiente sem \emph{heap} (embarcado,
\emph{kernel}), e para eliminar inicialização estática dinâmica.
\end{regra}

% =============================================================================
\subsection{Lambdas e objetos chamáveis}
\label{sec:lambdas}
% =============================================================================

O subprojeto \projeto{lambdas} parte de uma única ideia que explica praticamente
tudo o mais:

\begin{lstlisting}[style=bt, language=C++,
    caption={Uma \emph{lambda} é uma classe.}]
auto f = [x](int y) { return x + y; };

// o compilador gera, aproximadamente:
class __lambda {
    int x;                                        // a captura vira MEMBRO
public:
    auto operator()(int y) const { return x + y; }
};
\end{lstlisting}

Disso decorrem, sem nenhum acréscimo: por que \texttt{mutable} existe (o
\texttt{operator()} gerado é \texttt{const}); por que capturar por referência é
perigoso (o membro é uma referência, e ela pode pendurar); por que cada
\emph{lambda} tem tipo único (é uma classe anônima nova a cada expressão); e por
que \texttt{[=]} dentro de um método é a armadilha número um.

\subsubsection{Capturas}

\begin{lstlisting}[style=bt, language=bash,
    caption={Saída de \texttt{src/lambdas/}: o tamanho da classe gerada é a soma
             das capturas.}]
sizeof([]{})    = 1     ← classe VAZIA
sizeof([a]{})   = 4     = sizeof(int)
sizeof([&a]{})  = 8     = um ponteiro
sizeof([s]{})   = 32    = sizeof(std::string)
\end{lstlisting}

\begin{armadilha}[\texttt{[=]} num método captura \texttt{this} por ponteiro, não os membros]
\begin{lstlisting}[style=bt, language=C++, frame=none,
                   xleftmargin=0pt, xrightmargin=0pt]
[=] { return nome_; }
  // NÃO copia `nome_`
  // copia `this`, e `nome_` vira `this->nome_`
\end{lstlisting}
\begin{center}\footnotesize
\begin{tabular}{@{}lll@{}}
\toprule
\textbf{Captura} & \textbf{Copia} & \textbf{Seguro após o objeto morrer} \\
\midrule
\texttt{[this]}     & o ponteiro     & \textbf{não} \\
\texttt{[=]}        & o ponteiro (!) & \textbf{não} \\
\texttt{[*this]}    & o \textbf{objeto} & sim (C++17) \\
\texttt{[n=nome\_]} & só o membro    & sim, e é o mais barato \\
\bottomrule
\end{tabular}
\end{center}
O C++20 \textbf{depreciou} a captura implícita de \texttt{this} via
\texttt{[=]}. A regra: se a \emph{lambda} pode sobreviver ao objeto ---
\emph{callback}, \emph{thread}, \texttt{std::function} guardada --- nunca
capture \texttt{this}.
\end{armadilha}

A \emph{init-capture} do C++14 resolve quatro problemas com uma sintaxe só:
\texttt{[r = std::move(ptr)]} captura tipos \emph{move-only};
\texttt{[total = v.size() * 10]} pré-calcula uma vez em vez de a cada chamada;
\texttt{[h = cfg.host]} captura um \emph{membro} sem copiar o objeto inteiro; e
\texttt{[\&n = contador]} dá nome curto a algo de nome longo.

\begin{armadilha}[ordem de avaliação de argumentos, medida num bug real]
O próprio subprojeto encontrou o caso ao instrumentar uma \emph{lambda}
\texttt{mutable}: \texttt{print("\{\} \{\} \{\}", f(), f(), f())} tem ordem de
avaliação \textbf{não especificada}, e o GCC avalia da direita para a esquerda
--- imprimindo \texttt{3 2 1}. Quando o efeito colateral importa, use variáveis
nomeadas em vez de várias chamadas na mesma expressão.
\end{armadilha}

\subsubsection{Genéricas}

Uma \emph{lambda} com parâmetro \texttt{auto} é um \emph{template} de
\texttt{operator()}. A forma com \emph{template} explícito, do C++20, é
preferível sempre que o \emph{tipo} precisar de nome --- para
\texttt{std::forward}, para declarar variáveis daquele tipo, ou para acessar
\texttt{T::algo}:

\begin{lstlisting}[style=bt, language=C++,
    caption={A mesma \emph{lambda}, com e sem nome para o tipo.}]
[](auto&& x)          { return std::forward<decltype(x)>(x); }   // verboso
[]<typename T>(T&& x) { return std::forward<T>(x); }             // legível
\end{lstlisting}

E vale a mesma medição da \cref{sec:value-semantics}, agora dentro da
\emph{lambda}:

\begin{lstlisting}[style=bt, language=bash,
    caption={Saída de \texttt{src/lambdas/}: o parâmetro tem nome, logo é
             \emph{lvalue}.}]
sem forward, arg rvalue: 1 cópia(s), 0 move(s)     ← o parâmetro tem NOME, é lvalue
com forward, arg rvalue: 0 cópia(s), 1 move(s)
com forward, arg lvalue: 1 cópia(s), 0 move(s)     ← correto
\end{lstlisting}

\subsubsection{\texttt{std::function}: \emph{type erasure} e o seu preço}

\texttt{std::function} resolve o problema criado pelo fato de cada
\emph{lambda} ter tipo próprio: guardar chamáveis heterogêneos com a mesma
assinatura. O mecanismo é \gls{typeerasure}, e a estrutura é a mesma do
\texttt{std::any} (\cref{sec:vocabulary}) e do ECS (\cref{sec:ecs}):

\begin{lstlisting}[style=bt, language=C++,
    caption={O esqueleto do \emph{type erasure}.}]
struct Conceito { virtual R chamar(Args...) const = 0; };
template <typename F> struct Modelo : Conceito { F f_; /* ... */ };
std::unique_ptr<Conceito> impl_;
\end{lstlisting}

O tipo concreto some da assinatura mas continua vivo dentro do
\texttt{Modelo<F>} --- e é exatamente daí que vem o custo: a chamada passa por
uma \gls{vtable} e o corpo não pode sofrer \emph{inline}.

\begin{lstlisting}[style=bt, language=bash,
    caption={Saída de \texttt{src/lambdas/}: 2 milhões de elementos.}]
template<F>        :  12 ms   ← inlinável
std::function      :  43 ms   ← chamada virtual (3,5×)
ponteiro p/ função :  13 ms
\end{lstlisting}

A medição foi feita em \texttt{-O0}; em \texttt{-O2} a diferença
\textbf{aumenta}, porque o \emph{template} costuma desaparecer por completo
enquanto \texttt{std::function} mantém a indireção. Há ainda o custo de
alocação: a otimização de objeto pequeno cobre até certo tamanho de captura, e
acima disso cada \texttt{std::function} é um \emph{malloc}.

\begin{regra}[não use \texttt{std::function} como parâmetro de quem só chama e descarta]
É o erro mais comum da família. Uma função que recebe um comparador, o chama e o
esquece deve recebê-lo como \texttt{template <typename Comp>}: assim ele é
inlinado e não aloca. \texttt{std::function} é para \textbf{guardar} ---
membro, container, fila de tarefas --- e ali o custo se paga.
\end{regra}

\begin{table}[htbp]
\centering
\caption{As três formas de receber ou guardar um chamável.}
\label{tab:chamaveis-custo}
\footnotesize
\begin{tabular}{@{}lcccl@{}}
\toprule
\textbf{Ferramenta} & \textbf{Inlina} & \textbf{Aloca} & \textbf{Guardável} & \textbf{Assinatura} \\
\midrule
\texttt{template <typename F>} & \textbf{sim} & não    & não          & genérica \\
ponteiro para função           & não          & não    & sim          & fixa \\
\texttt{std::function}         & não          & talvez & \textbf{sim} & fixa \\
\bottomrule
\end{tabular}
\end{table}

\subsubsection{\classe{std::invoke} e ponteiros para membro}

A linguagem trata \textbf{cinco} sintaxes de chamada de forma diferente ---
\texttt{f(args)}, \texttt{(obj.*pm)(args)}, \texttt{(ptr->*pm)(args)},
\texttt{obj.*pd}, \texttt{ptr->*pd} --- e escrever código genérico sobre elas
exigiria um \texttt{if constexpr} para cada caso. \texttt{std::invoke} unifica
as cinco, e é por isso que ele aparece na implementação de \texttt{std::thread},
\texttt{std::async} e de qualquer \emph{thread pool} (\cref{sec:concurrency}).

O uso mais prático de um ponteiro para membro, no entanto, é como
\textbf{projeção} nos algoritmos de \emph{ranges}:
\texttt{std::ranges::sort(equipe, \{\}, \&Pessoa::idade)}
(\cref{sec:ranges}).

\begin{regra}[\emph{lambda} em vez de \texttt{std::bind}]
\texttt{std::bind} usa \emph{placeholders} opacos, copia os argumentos ligados
por padrão, gera erros ilegíveis quando aninhado e trata \texttt{bind} dentro de
\texttt{bind} de forma surpreendente. As \emph{Core
Guidelines}~\cite{cppcoreguidelines} recomendam \emph{lambda} (F.50). Quando a
aplicação parcial é literalmente ``fixar os primeiros argumentos'',
\classe{std::bind\_front} (C++20) é legível, inlinável e não aloca.
\end{regra}

\subsubsection{Cinco padrões que só existem por causa das \emph{lambdas}}

\paragraph{IIFE.} Inicializar um \texttt{const} que precisa de várias instruções,
sem criar uma função auxiliar --- e sem abrir mão do \texttt{const}:

\begin{lstlisting}[style=bt, language=C++,
    caption={\emph{Immediately Invoked Function Expression}.}]
const std::string prefixo = [nivel] {
    if (nivel >= 3) return "[DEBUG] ";
    if (nivel >= 2) return "[INFO] ";
    return "[ERRO] ";
}();
\end{lstlisting}

É praticamente obrigatório para inicializar um \emph{membro} \texttt{const} com
lógica na lista de inicialização (\cref{sec:syntax}).

\paragraph{\emph{Overload idiom}.} Já apresentado na \cref{sec:vocabulary}; o
que o subprojeto acrescenta é a explicação de por que ele funciona --- e por que
a terceira linha é indispensável. \texttt{Overload : Ts...} herda de todas as
\emph{lambdas}, mas sem \texttt{using Ts::operator()...;} os
\texttt{operator()} ficariam escondidos uns pelos outros: é o mesmo \emph{name
hiding} da \cref{sec:oop}.

\paragraph{Recursão.} Uma \emph{lambda} não tem nome dentro do próprio corpo. As
três saídas diferem em custo e em disponibilidade
(\cref{tab:recursao-lambda}); enquanto o \emph{deducing this} não chega ao GCC
13, o \emph{y-combinator} é a opção certa em código quente.

\begin{table}[htbp]
\centering
\caption{Recursão em \emph{lambda}: as três formas.}
\label{tab:recursao-lambda}
\footnotesize
\begin{tabular}{@{}lcccc@{}}
\toprule
\textbf{Forma} & \textbf{Aloca} & \textbf{Inlina} & \textbf{Chamada} & \textbf{GCC 13} \\
\midrule
\texttt{std::function} & talvez & não & \texttt{f(n)}    & sim \\
\emph{y-combinator}    & não    & sim & \texttt{f(f, n)} & sim \\
\emph{deducing this}   & não    & sim & \texttt{f(n)}    & \textbf{não} \\
\bottomrule
\end{tabular}
\end{table}

\paragraph{\emph{Scope guard}.} \gls{raii} para \emph{qualquer} ação, escrito no
ponto de uso, sem uma classe por recurso. É a generalização do que a
\cref{sec:smart-pointers} fazia com \emph{deleters} customizados --- e, com o
\texttt{dispensar()}, cobre também o caso ``desfaz, a menos que dê certo'':

\begin{lstlisting}[style=bt, language=C++,
    caption={Limpeza e \emph{rollback} com o mesmo mecanismo.}]
auto g = GuardaDeEscopo{[&]{ fclose(f); }};
auto desfazer = GuardaDeEscopo{[&]{ rollback(); }};
// ... deu certo:
desfazer.dispensar();                      // commit
\end{lstlisting}

\paragraph{\emph{Lambda} como comparador de container.} Desde o C++20, uma
\emph{lambda} sem captura é construtível por padrão, o que permite usá-la
diretamente como argumento de \emph{template}:
\texttt{std::set<std::string, PorTamanho>}. A consequência a lembrar é que o
comparador faz parte do \textbf{tipo}: dois \emph{sets} com \emph{lambdas}
distintas são incompatíveis mesmo tendo a mesma lógica.

% =============================================================================
\subsection{Síntese}
% =============================================================================

Os quatro subprojetos deste capítulo são quatro respostas à mesma pergunta:
\emph{quanto do programa dá para decidir antes de rodar?}

\emph{Templates} decidem o \emph{código}, gerando uma versão especializada por
tipo. \emph{Concepts} decidem a \emph{admissibilidade}, movendo o erro para o
ponto da chamada. \texttt{constexpr} decide os \emph{valores}, e leva junto uma
garantia que nenhum teste dá: naquele caminho não há \gls{ub}. E as
\emph{lambdas} decidem o \emph{chamável}, com a ressalva de que
\texttt{std::function} desfaz boa parte do ganho ao trocar o tipo estático por
uma indireção.

\begin{regra}[a hierarquia de escolha em código genérico novo]
Na ordem: (1)~se dá para restringir com um \emph{concept}, use \emph{concept};
(2)~se dá para escolher o galho com \texttt{if constexpr}, use \texttt{if
constexpr}; (3)~SFINAE apenas para manter código antigo ou para atingir
compiladores anteriores ao C++20. E, para chamáveis: \texttt{template <typename
F>} por padrão, \texttt{std::function} só quando for preciso \emph{guardar}.
\end{regra}
